% =====================================================================
%  LAPORAN TEKNIS — PhishGuard AI
%  Template LaTeX. Kompilasi: pdflatex laporan-teknis.tex (2x untuk ToC)
%
%  Bagian yang PERLU DIISI ditandai  % >>> ISI:  — cari dengan Ctrl-F.
%  Angka teknis di dalam sudah diisi dari pengukuran nyata proyek;
%  jangan diubah tanpa memverifikasi ulang ke kode/APK.
% =====================================================================
\documentclass[11pt,a4paper]{article}

% ── Paket ────────────────────────────────────────────────────────────
\usepackage[T1]{fontenc}
\usepackage[utf8]{inputenc}
% babel sengaja tidak dipakai: locale Indonesia tidak ada di semua
% instalasi TeX. Teks Indonesia tetap render sempurna tanpa babel.
\usepackage[a4paper,margin=2.5cm]{geometry}
\usepackage{graphicx}
\usepackage{booktabs}
\usepackage{tabularx}
\usepackage{xcolor}
\usepackage{listings}
\usepackage{enumitem}
\usepackage{fancyhdr}
\usepackage{hyperref}
\usepackage{caption}

% ── Warna brand ──────────────────────────────────────────────────────
\definecolor{brandblue}{HTML}{2563EB}
\definecolor{navy}{HTML}{0F172A}
\definecolor{slate}{HTML}{64748B}
\definecolor{codebg}{HTML}{F1F5F9}

% ── Hyperlink ────────────────────────────────────────────────────────
\hypersetup{
  colorlinks=true, linkcolor=brandblue, urlcolor=brandblue,
  citecolor=brandblue, pdftitle={Laporan Teknis PhishGuard AI},
}

% ── Judul bagian (tanpa paket tambahan; hanya xcolor) ────────────────
\makeatletter
\renewcommand\section{\@startsection{section}{1}{\z@}%
  {-3.5ex \@plus -1ex \@minus -.2ex}{1.3ex \@plus.2ex}%
  {\normalfont\Large\bfseries\color{navy}}}
\renewcommand\subsection{\@startsection{subsection}{2}{\z@}%
  {-3.25ex\@plus -1ex \@minus -.2ex}{1ex \@plus.2ex}%
  {\normalfont\large\bfseries\color{brandblue}}}
\makeatother

% ── Header/footer ────────────────────────────────────────────────────
\pagestyle{fancy}
\fancyhf{}
\lhead{\small\color{slate}PhishGuard AI}
\rhead{\small\color{slate}Laporan Teknis}
\cfoot{\thepage}
\renewcommand{\headrulewidth}{0.4pt}

% ── Gaya kode ────────────────────────────────────────────────────────
\lstset{
  backgroundcolor=\color{codebg},
  basicstyle=\ttfamily\small,
  keywordstyle=\color{brandblue}\bfseries,
  commentstyle=\color{slate}\itshape,
  stringstyle=\color{navy},
  numbers=none, frame=single, rulecolor=\color{slate!40},
  breaklines=true, showstringspaces=false,
  aboveskip=1em, belowskip=1em,
}

% ── Kotak ringkas ────────────────────────────────────────────────────
\newcommand{\keyfact}[2]{\textbf{\color{brandblue}#1}\,: #2}

% =====================================================================
\begin{document}

% ── HALAMAN JUDUL ────────────────────────────────────────────────────
\begin{titlepage}
  \centering
  \vspace*{2cm}

  % >>> ISI: ganti dengan logo bila mau — taruh file di folder yang sama
  % \includegraphics[width=0.35\textwidth]{logo.png}\\[1.5cm]

  {\Huge\bfseries\color{navy} PhishGuard AI\par}
  \vspace{0.4cm}
  {\Large\color{slate} Sistem Perlindungan Anti-Penipuan Digital\\
   Berpusat pada Manusia untuk Pengguna Indonesia\par}
  \vspace{2cm}

  {\large\bfseries LAPORAN TEKNIS\par}
  \vspace{0.3cm}
  {\color{slate} Versi Aplikasi 2.0\par}
  \vspace{3cm}

  \begin{tabular}{rl}
    \textbf{Tim} & \rule{5cm}{0.4pt} \\[4pt]         % >>> ISI: nama tim
    \textbf{Anggota} & \rule{5cm}{0.4pt} \\[4pt]     % >>> ISI: anggota 1
                     & \rule{5cm}{0.4pt} \\[4pt]     % >>> ISI: anggota 2
    \textbf{Track} & Human-Centric Security \\[4pt]
    \textbf{Tanggal} & \rule{5cm}{0.4pt} \\[4pt]     % >>> ISI: tanggal
    \textbf{Repositori} & \url{https://github.com/fsrsr/Phisguard-AI} \\
  \end{tabular}

  \vfill
  {\small\color{slate} Dokumen ini dibuat sebagai template; bagian bertanda
  ``ISI'' perlu dilengkapi. Seluruh angka teknis berasal dari pengukuran
  langsung terhadap kode dan APK.\par}
\end{titlepage}

\tableofcontents
\newpage

% =====================================================================
\section{Ringkasan Eksekutif}
% =====================================================================
PhishGuard AI adalah aplikasi Android yang melindungi pengguna Indonesia
dari penipuan digital pada \emph{titik paparan}: tautan phishing di
peramban, pesan penipuan di SMS dan WhatsApp, situs judi online (judol),
dan panggilan penipuan. Berbeda dari alat keamanan konvensional, setiap
peringatan dijelaskan dalam bahasa Indonesia yang dapat dipahami, sehingga
tiap ancaman menjadi kesempatan belajar bagi pengguna.

Prinsip yang mengikat seluruh keputusan teknisnya:
\begin{quote}
  \itshape ``Peringatan palsu sama merusaknya dengan ancaman yang
  terlewat.''
\end{quote}
Alat yang terlalu sering salah akan diabaikan, dan sejak itu berhenti
melindungi. Karena itu penekanan utama proyek ini bukan sekadar mendeteksi
lebih banyak, melainkan \textbf{menekan salah-blokir} sambil tetap
menangkap ancaman nyata.

\vspace{0.5em}
\noindent\textbf{Capaian terukur:}
\begin{itemize}[leftmargin=1.4em,itemsep=2pt]
  \item Salah-blokir domain populer turun dari \textbf{6{,}02\%} menjadi
        \textbf{0{,}45\%}.
  \item \textbf{4{,}1 juta} domain ancaman diperiksa luring dalam aset
        \textbf{9{,}4\,MB}.
  \item Kosakata judol dikalibrasi dari \textbf{1{,}8 juta} domain judi
        nyata.
  \item Terverifikasi berjalan di perangkat nyata (Vivo V2322, Android~14).
\end{itemize}

% =====================================================================
\section{Latar Belakang dan Rumusan Masalah}
% =====================================================================
Studi Kaspersky (2023) menemukan 64\% insiden siber berakar pada kesalahan
manusia. Di Indonesia, paparannya akut: phishing perbankan, tautan judol
yang menyebar lewat WhatsApp, serta SMS penipuan yang mengatasnamakan
kurir, bantuan sosial, dan pinjaman ilegal.

Perlindungan yang ada gagal melayani pengguna Indonesia dalam tiga hal:
\begin{enumerate}[leftmargin=1.6em,itemsep=3pt]
  \item \textbf{Berbahasa Inggris} — peringatan tidak terbaca oleh yang
        paling rentan.
  \item \textbf{Tidak tersetel untuk ancaman lokal} — judol terlewat sama
        sekali, begitu pula spam yang sengaja dikaburkan.
  \item \textbf{Terlalu sering salah} — pemblokiran berlebih membuat
        pengguna belajar bahwa peringatan biasanya keliru, lalu
        mengabaikannya secara refleks, termasuk yang benar.
\end{enumerate}

Pemblokiran berlebih adalah mode kegagalan yang jarang diukur. Pengujian
pendekatan berbasis kata kunci konvensional menunjukkan ia menandai 6\%
domain terpopuler dunia sebagai phishing, dan memblokir setiap situs yang
memiliki halaman \texttt{/login}. Selain itu, peringatan yang hanya
berbunyi ``PHISHING DETECTED'' tidak mengajarkan apa pun; pengguna tidak
memperoleh kemampuan mengenali serangan berikutnya.

% =====================================================================
\section{Arsitektur Sistem}
% =====================================================================
Aplikasi ditulis dalam Kotlin (\textasciitilde8.700 baris),
\keyfact{minSdk}{24 (Android 7.0)}, \keyfact{targetSdk}{36}. Seluruh
analisis berjalan di perangkat; tidak ada data penjelajahan yang dikirim
ke luar.

\subsection{Empat Kanal Masuk}
\begin{center}
\begin{tabularx}{\textwidth}{@{}l X l@{}}
  \toprule
  \textbf{Kanal} & \textbf{Yang dipantau} & \textbf{Mekanisme Android} \\
  \midrule
  Peramban & URL di 24 peramban dan WebView & \texttt{AccessibilityService} \\
  SMS & Pesan masuk aplikasi bawaan & \texttt{BroadcastReceiver} \\
  Chat \& notifikasi & WhatsApp, Telegram, dll. & \texttt{NotificationListenerService} \\
  Panggilan & Transkripsi percakapan langsung & \texttt{SpeechRecognizer} + FGS \\
  \bottomrule
\end{tabularx}
\end{center}

\subsection{Pipeline Analisis URL}
\texttt{AnalysisPipeline} memeriksa berurutan, dari yang paling pasti ke
yang paling menebak. Pemeriksaan yang dapat keliru selalu ditaruh
\emph{setelah} daftar putih, sehingga kekeliruannya tidak pernah mengenai
situs tepercaya.

\begin{center}
\begin{tabularx}{\textwidth}{@{}c X l@{}}
  \toprule
  \# & \textbf{Langkah} & \textbf{Bila cocok} \\
  \midrule
  1 & Daftar putih domain resmi (BCA, BRI, OJK, \ldots) & \textbf{aman, berhenti} \\
  2 & Database phishing lokal (URL persis) & ancaman 1{,}00 \\
  3 & Database phishing per-host & ancaman $\geq$0{,}90 \\
  4 & Blocklist global (Bloom filter, 4{,}1 juta domain) & ancaman 0{,}95 \\
  5 & Skoring heuristik (\texttt{UrlRiskScorer}) & sesuai skor \\
  \bottomrule
\end{tabularx}
\end{center}

Daftar putih sengaja didahulukan: entri dari laporan komunitas maupun umpan
publik bisa keliru atau disuntik pihak lain, dan domain bank tidak boleh
dapat dijatuhkan oleh sumber semacam itu.

% =====================================================================
\section{Solusi: Empat Keputusan Desain}
% =====================================================================

\subsection{Skoring Berbobot, Bukan Pencocokan Kata Kunci}
Setiap sinyal hanya \emph{menyumbang} bobot; hanya akumulasinya yang
bertindak. Satu sinyal lemah tidak pernah cukup untuk memblokir.

\begin{center}
\begin{tabularx}{\textwidth}{@{}X c X@{}}
  \toprule
  \textbf{Sinyal} & \textbf{Bobot} & \textbf{Catatan} \\
  \midrule
  Peniruan brand & 0{,}75 & 0{,}45 untuk brand yang juga kata umum \\
  Istilah judi pada hostname & 0{,}72 & tidak memblokir sendirian \\
  Alamat IP mentah & 0{,}55 & \\
  Domain punycode (\texttt{xn-{}-}) & 0{,}50 & huruf menyerupai latin \\
  TLD terbukti disalahgunakan & 0{,}30 & dikalibrasi dari data \\
  Kata sensitif di \emph{path} & 0{,}12 & hampir semua situs punya \texttt{/login} \\
  \bottomrule
\end{tabularx}
\end{center}

Dampaknya terukur: kata \texttt{login} di path yang dulu memicu blokir
0{,}88 kini hanya 0{,}12. Pencocokan brand tidak pernah memakai substring
bebas --- hanya token utuh, token+afiks dikenal (\texttt{klikbca}), atau
typosquat berjarak edit 1 --- sehingga \texttt{bri-mobile.xyz} tertangkap
tetapi \texttt{bridge.com} tidak.

\subsection{Respons Bertingkat}
Peringatan yang salah tidak boleh menjegal pengguna.
\begin{center}
\begin{tabularx}{\textwidth}{@{}l X@{}}
  \toprule
  \textbf{Skor} & \textbf{Tindakan} \\
  \midrule
  $\geq$ 0{,}80 & Layar peringatan penuh; situs dihentikan \\
  0{,}55--0{,}80 & Notifikasi yang bisa diabaikan; penjelajahan lanjut \\
  $<$ 0{,}55 & Hanya dicatat di log \\
  \bottomrule
\end{tabularx}
\end{center}
Deteksi yang belum pasti tidak membajak layar, sehingga pengguna tetap
berpikir alih-alih patuh membabi buta --- sekaligus menjawab kekhawatiran
bahwa otomasi menurunkan kewaspadaan.

\subsection{Alasan Berbahasa Manusia}
Alih-alih ``PHISHING DETECTED'', pengguna membaca \emph{mengapa}:
\begin{quote}
  \itshape ``Memakai nama `bca' padahal domainnya (klikbca-verifikasi.tk)
  bukan milik resminya. Nama domain memuat kata `verifikasi' yang khas
  halaman pencuri data. Domain memakai akhiran .tk yang banyak dipakai
  penipuan.''
\end{quote}
Pengguna mempelajari konsep peniruan brand tepat saat mengalaminya ---
bentuk paling efektif dalam pendidikan keamanan.

\subsection{Human-in-the-Loop}
Pengguna dapat membantah vonis sistem dalam satu ketukan, dan keputusan itu
diingat permanen lintas sesi. Sistem boleh salah; pengguna yang memutuskan.

% =====================================================================
\section{Komponen Kunci}
% =====================================================================

\subsection{Blocklist Luring 4,1 Juta Domain}
Sumber: blocklist publik Lyra (ITSEC Research). Data mentah lebih dari
80\,MB dipampatkan menjadi \textbf{Bloom filter 9{,}4\,MB} berisi
4.108.088 domain (phishing, malware, judi), dipetakan ke memori
(\texttt{mmap}) langsung dari APK dengan pencarian $O(1)$ tanpa jaringan.

Sifat Bloom filter: dapat keliru menyatakan sebuah domain terdaftar, tapi
tidak pernah sebaliknya. Peluangnya disetel 0{,}01\% dan terukur 0{,}0085\%.
Domain resmi diperiksa lebih dulu, sehingga kekeliruan itu tidak mengenai
situs tepercaya.

\subsection{Kalibrasi dari Data, Bukan Tebakan}
TLD tersangka dan kata penanda judi \emph{tidak} ditebak. Dengan
membandingkan 4{,}1 juta domain jahat terhadap 1 juta domain populer
(Tranco), dihitung \emph{lift} tiap TLD --- pangsanya di daftar jahat
dibagi pangsanya di daftar sah:
\begin{center}
\begin{tabular}{@{}lr lr@{}}
  \toprule
  TLD & lift & TLD & lift \\
  \midrule
  \texttt{.tk} & 88$\times$ & \texttt{.click} & 21$\times$ \\
  \texttt{.gdn} & 86$\times$ & \texttt{.sbs} & 18$\times$ \\
  \texttt{.stream} & 82$\times$ & \texttt{.cyou} & 13$\times$ \\
  \bottomrule
\end{tabular}
\end{center}
Kata judol \texttt{rtp} (lift 177$\times$), \texttt{gacor} (86$\times$),
dan \texttt{mpo} (57$\times$) diadopsi. Kata Indonesia umum seperti
\texttt{daftar} dan \texttt{situs} sengaja ditolak meski lift-nya tinggi,
karena acuan Tranco kurang memuat situs berbahasa Indonesia.

\subsection{Deteksi Spam yang Dikaburkan}
Spam judol lokal ditulis untuk lolos penyaring kata:
\begin{center}
\ttfamily D e p 0 \quad M4X'W1N \quad Mahy'0ng \quad Gratiss
\end{center}
Modul \texttt{TextNormalizer} membalik empat teknik sekaligus (huruf
berspasi, leetspeak, tanda baca penyisip, huruf berulang) sehingga
terurai menjadi \texttt{depo}, \texttt{maxwin}, \texttt{mahjong},
\texttt{gratis}. Hasil normalisasi dipakai \emph{bersama} teks asli, bukan
menggantikannya --- melipat angka merusak makna sah seperti ``100 juta''.

\subsection{Peta Pola Ancaman}
Fitur opsional yang memetakan ancaman terkonfirmasi ke tiga dimensi:
spasial (titik rawan), temporal (rentang waktu), dan kategorikal, dengan
ekspor GeoJSON. Dibangun dengan tiga pagar privasi: \emph{opt-in} dan mati
secara bawaan, lokasi kasar yang dibulatkan \textasciitilde1\,km, dan
tidak pernah melacak --- hanya membaca lokasi terakhir yang diketahui
sistem saat ancaman dikonfirmasi.

% =====================================================================
\section{Keamanan dan Privasi}
% =====================================================================
\begin{itemize}[leftmargin=1.4em,itemsep=3pt]
  \item \textbf{Analisis di perangkat.} Yang keluar hanya URL yang
        \emph{dikonfirmasi pengguna} sebagai penipuan, untuk basis data
        komunitas. Riwayat penjelajahan tidak pernah dikirim.
  \item \textbf{Aturan Firestore.} Membatasi bentuk tulisan, mengunci
        daftar putih resmi dari klien, dan menutup penghapusan --- sudah
        diterapkan dan diverifikasi terhadap basis data produksi.
  \item \textbf{Isi pesan tidak dicatat.} Notifikasi chat adalah
        percakapan pribadi; hanya panjangnya yang dicatat.
  \item \textbf{Integritas aplikasi.} \texttt{AppIntegrityChecker}
        mencocokkan sidik jari sertifikat penandatangan; APK yang
        ditandatangani ulang menolak berjalan.
\end{itemize}

% =====================================================================
\section{Pengujian dan Verifikasi}
% =====================================================================
Logika deteksi diuji oleh \textbf{46 unit test} murni JVM (tanpa
emulator), yang mengunci perilaku dua arah: situs normal tidak boleh
diblokir, dan pola phishing nyata harus tertangkap.

\begin{center}
\begin{tabularx}{\textwidth}{@{}l c X@{}}
  \toprule
  \textbf{Berkas uji} & \textbf{Kasus} & \textbf{Fokus} \\
  \midrule
  \texttt{UrlRiskScorerTest} & 20 & FP situs normal + deteksi phishing nyata \\
  \texttt{TextNormalizerTest} & 11 & spam yang dikaburkan \\
  \texttt{MessageRuleEngineTest} & 10 & pola pesan SMS/chat \\
  \texttt{DomainBlocklistTest} & 4 & kesepakatan algoritma Python$\leftrightarrow$Kotlin \\
  \bottomrule
\end{tabularx}
\end{center}

Verifikasi lapangan dilakukan pada perangkat nyata (Vivo V2322,
Android~14): deteksi peramban, judol berbasis isi halaman, respons
bertingkat, navigasi kembali setelah blokir, dan intersepsi notifikasi
WhatsApp semuanya dikonfirmasi lewat log dan tangkapan layar.

% =====================================================================
\section{Batasan dan Pengembangan Lanjut}
% =====================================================================
Laporan ini menyatakan batasan secara terbuka:
\begin{itemize}[leftmargin=1.4em,itemsep=3pt]
  \item \textbf{Berbasis heuristik, bukan model pembelajaran mesin.}
        Deteksi memakai skoring berbobot dan aturan yang dikalibrasi dari
        data nyata --- pendekatan yang justru menghasilkan angka
        false-positive yang dapat dibuktikan. Integrasi model TFLite
        disiapkan sebagai jalur pengembangan.
  \item \textbf{Bergantung pada bilah alamat peramban} lewat
        \texttt{viewId}; peramban yang tidak dikenal dapat lolos.
  \item \textbf{Ambang belum tervalidasi lapangan} dalam skala besar;
        angka 0{,}55 dan 0{,}80 dikalibrasi dari contoh.
\end{itemize}

Arah lanjut: mode latihan interaktif dari data ancaman nyata pengguna
(kuis ``mana yang asli?''), umpan balik perilaku bulanan, dan validasi
ambang dari data lapangan.

% =====================================================================
\section*{Lampiran: Reproduksi Angka}
\addcontentsline{toc}{section}{Lampiran: Reproduksi Angka}
% =====================================================================
Seluruh angka dapat direproduksi dari kode:
\begin{lstlisting}[language=bash]
# False positive & cakupan blocklist
python3 tools/build_blocklist.py --lyra <Lyra> --tranco top-1m.csv.zip
# -> tools/blocklist_report.md

# Seluruh unit test
./gradlew testDebugUnitTest

# APK rilis
./gradlew assembleRelease
\end{lstlisting}

\end{document}
